%--------------------------------------------------------------------------------
%
%
%--------------------------------------------------------------------------------
\pagebreak
\section*{FxCA - flush from data cache}
\addcontentsline{toc}{section}{FxCA - flush from cache}

\begin{iPageDescEntry}{Syntax}
    \texttt{FICA RegR, [ RegX ] ( RegB )} \\
    \texttt{FDCA RegR, [ RegX ] ( RegB )}
\end{iPageDescEntry}

\begin{iPageDescEntry}{Format}
	The \texttt{FICA} instruction uses the following instruction format. \\

    \begin{flushleft}
    \drawInstrFormat
        {0,1,3,5,5.5,6,6.5,8.5,9.5,11.5,16}
        {\OpCodeGrpSYS, \OpCodePCA, RegR, 0, M, 0, RegB, 3, RegX, 0}
    \end{flushleft}
    
    The \texttt{FDCA} instruction uses the following instruction format. \\

    \begin{flushleft}
    \drawInstrFormat
        {0,1,3,5,5.5,6,6.5,8.5,9.5,11.5,16}
        {\OpCodeGrpSYS, \OpCodePCA, RegR, 0, M, 1, RegB, 3, RegX, 0}
    \end{flushleft}
     
\end{iPageDescEntry}

\begin{iPageDescEntry}{Description}
    The \texttt{FCA} instruction flushes a cache line, selected by the effective address in \texttt{RegB}. If the M bit is set, the base register RegB is updated with the computed effective address. This is a privileged instruction.
\end{iPageDescEntry}

\begin{iPageDescOpEntry}{Operation}
\begin{lstlisting}[style=iPageOpStyle]
RegR <- flushFromCache( addOfs32( RegB, RegX ));
\end{lstlisting}
\end{iPageDescOpEntry}

\begin{iPageDescEntry}{Exceptions}
    Privilege violation trap.
\end{iPageDescEntry}

\begin{iPageDescEntry}{Notes}
None.
\end{iPageDescEntry} 
